\section{토의 및 한계점}
\label{sec:토의}

\subsection{제안 모델의 장점과 실용적 가치}
본 논문에서 제안한 \method\ 프레임워크는 전역 정보를 온보드 통신 오버헤드 없이 활용할 수 있는 글로벌-로컬 정책 증류 기법을 채택하였다. 오프라인에서 풍부하게 훈련된 교사의 지식을 경량 MLP 형태로 내재화하고, 스위치 게이트(PRS)를 통해 위험 시에만 복잡한 예측 연산을 선택적으로 호출함으로써, UAV 내부의 계산 지연을 줄이고 실시간 데이터 전달 속도를 보장하였다. 이는 전력 및 계산 능력이 제한된 실제 UAV swarms 통신 시스템에 바로 융합될 수 있는 강한 배포 실용성을 가진다.

\subsection{추가 검증의 필요성 및 한계점}
본 연구의 결과를 완벽한 투고 수준으로 확보하기 위해 극복해야 할 세부 한계점은 다음과 같다.
\begin{enumerate}
    \item \textbf{Phase 13 P+의 통계적 정밀화}: P+ 확장 모듈의 성능 기여는 프로토타입 시뮬레이션을 통해 동작이 입증되었으나, 최종 논문 게재를 위해서는 14개 전체 시나리오에 대한 다중 시드 독립 실행과 신뢰구간(Confidence Interval) 검증 데이터가 완전히 기재되어야 한다.
    \item \textbf{물리 계층 모사도의 정교화}: 현재 Python 기반 시뮬레이터(Lite-GLOBE) 환경은 전송 큐와 링크 지속 시간을 반영하지만, 다중 경로 간섭(Multipath fading), 노드의 순간 회전에 따른 안테나 지향성 변화, MAC 계층의 패킷 충돌 등 정밀한 물리/MAC 레벨의 영향이 생략되어 있다. 이를 보완하기 위해 ns-3 또는 OPNET과 같은 공인 통신 시뮬레이터를 활용한 교차 검증이 요구된다.
    \item \textbf{전통 레거시 라우팅 Baselines 추가}: 비교군에 GPSR, Evo-QGeo 외에 실제 네트워크 환경에서 널리 쓰이는 AODV, OLSR과 같은 전통적인 애드혹 프로토콜을 탑재하여, 홉 마다 누적되는 총 제어 오버헤드 바이트를 비교 분석할 필요가 있다.
    \item \textbf{통신 제어 오버헤드의 정량화}: 입력 풋프린트 바이트 외에, 패킷 전송 시 실시간으로 수반되는 비콘 주기와 헬로 메시지의 헤더 바이트 크기를 총합하여 종단간 통신 오버헤드 메트릭을 고도화해야 한다.
\end{enumerate}
